\documentclass[12pt,a4paper]{article}
\usepackage[utf8]{inputenc}
\usepackage[T1]{fontenc}
\usepackage{lmodern}
\usepackage{microtype}
\usepackage{amsmath, amssymb}
\usepackage{geometry}
\usepackage{graphicx}
\usepackage{setspace}
\usepackage{cite}
\usepackage{titlesec}
\usepackage[colorlinks=true, linkcolor=blue, urlcolor=blue, citecolor=black]{hyperref}

\geometry{margin=2.5cm}
\setstretch{1.3} % Line spacing improving readability
\titleformat{\section}{\bfseries\Large}{\thesection.}{1em}{}
\titlespacing*{\section}{0pt}{1.5ex plus 0.5ex minus 0.2ex}{0.8ex plus 0.2ex}
\titleformat{\subsection}{\bfseries\normalsize}{\thesubsection.}{1em}{}
\titlespacing*{\subsection}{0pt}{1.2ex plus 0.4ex minus 0.2ex}{0.5ex}

\title{\textbf{Stationary Strings Theorem: Relativity, Phase Transition and the Cosmological Cycle of Information}\\
\large (The Bridge Between Relativity and Stationary Strings -- Extended Version)}
\author{Fahri Meral}
\date{2025}

\begin{document}

\maketitle

\begin{abstract}
This study aims to reinterpret the universe through two fundamental states of existence: the \textit{vibrating energy phase} and the \textit{stationary information phase}.
Einstein's general theory of relativity predicts that time slows down in strong gravitational fields and stops at the event horizon.
String theory, on the other hand, asserts that vibrating energy strings lie at the foundation of the universe.
The \textbf{Stationary Strings Theorem}, which unites these two perspectives, ties the flow of time to the vibration frequency ($f$) of strings.

In this study, a new physical quantity called ``Perfection Potential'' ($\Phi$) is proposed as a measure of order. This potential is defined as the inverse of entropy ($S$): $\Phi = 1/S(f)$. Black holes are cosmic mechanisms where string vibration halts ($f \to 0$) and entropy approaches zero ($S \to 0$); in that case the Perfection Potential diverges to infinity ($\Phi \to \infty$). This phase transition shows that energy does not vanish, but is transformed into stationary information ruled by absolute order. The universe thus exists in a continuous cycle between information and energy.
\end{abstract}

\section{Introduction}
With his general theory of relativity, Albert Einstein showed that time is not absolute, but a dimension curved by the gravitational field.
Similarly, string theory proposes that all particles in the universe consist of strings vibrating at different frequencies.
At the foundation of both approaches lies the \textit{time--energy relationship}.

Combining these two frameworks, this article puts forward the following fundamental proposition:
\begin{quote}
``The slowing of time is a direct consequence of the decrease in the vibration frequency of strings.
The stopping of time is the transition of strings into the stationary state, where they lose their vibration and perfect order and pure information reign.''
\end{quote}

Under this assumption, black holes perform a \textit{phase transition} in the universe:
the vibrating energy form passes into the stationary information form.
This information is the existential trace of energy; that is, the ``informational DNA'' of the universe.

\section{Time Dilation in Relativity}

In general relativity, time is directly related to the gravitational potential.
As an observer approaches a black hole, the flow of time slows down and stops completely at the event horizon.

The time dilation equation is:
\begin{equation}
t' = t \sqrt{1 - \frac{2GM}{r c^2}}
\end{equation}

Here:
\begin{itemize}
\item $t'$: the time measured by the local observer,
\item $t$: the time measured by the distant observer,
\item $G$: the universal gravitational constant,
\item $M$: the mass of the black hole,
\item $r$: the distance of the observer from the center of the black hole,
\item $c$: the speed of light.
\end{itemize}

As $r \to \frac{2GM}{c^2}$, the inside of the root approaches zero and $t' \to 0$.
This means time stops at the event horizon.

\subsection{The Event Horizon and the Slowing of Time}
The event horizon is not merely a boundary from which light cannot escape, but a point where time loses its meaning.
Beyond this boundary the classical concept of time collapses; this can be directly associated with the slowing of string vibration.

\section{The Relation Between String Frequency and Time}
In string theory, energy is directly proportional to the frequency of the string:
\begin{equation}
E = h f
\end{equation}
and together with the Einstein mass-energy equivalence:
\begin{equation}
E = mc^2 \Rightarrow m = \frac{h f}{c^2}
\end{equation}
Hence the energy of a string decreases as its vibration frequency decreases.

Under gravity, frequency changes in a manner analogous to time dilation:
\begin{equation}
f' = f_0 \sqrt{1 - \frac{2GM}{r c^2}}
\end{equation}
At the event horizon $f' \to 0$; that is, the string loses its vibration.

\subsection{Variation of Frequency with Gravity}
This expression shows that a direct link can be established between time and frequency:
\begin{equation}
t \propto f
\end{equation}
Consequently, the approach of frequency to zero leads to the stopping of time.

\section{The Unified Equation and the Stationarity of Time}
The general formula combining time and string frequency:
\begin{equation}
t' = k_t \cdot f_0 \sqrt{1 - \frac{2GM}{r c^2}}
\end{equation}
Here the constant $k_t$ represents the perceptual scale of time.
As $r \to \frac{2GM}{c^2}$, $t' \to 0$ and $f' \to 0$;
that is, string vibration halts and time loses its meaning.

\section{The Two Phases of the Universe: Energy and Information}

\subsection{The Vibrating Energy Phase}
This phase is the observable universe: strings vibrate ($f > 0$), energy and matter interact, entropy increases.
The arrow of time points forward.
In this phase order continuously degrades and the universe reaches ever higher entropy.

\subsection{The Stationary Information Phase (Perfection Dimension)}
In this phase the vibration of strings halts ($f = 0$).
The concepts of energy, time and entropy lose their meaning.
This is the ``perfect state'' into which matter transforms.
Information becomes absolute order --- there is a single microstate of the system.

\subsection{Perfection Potential ($\Phi$) and Mathematical Model}
To measure order and chaos in the universe, we may define a new physical quantity called the ``Perfection Potential'' ($\Phi$). This potential is the inverse of the entropy function ($S(f)$) depending on the string frequency of the system:
\begin{equation}
\Phi = \frac{1}{S(f)}
\end{equation}
By this definition:
\begin{itemize}
    \item As vibration increases ($f > 0$), entropy increases ($S$ increases) and the Perfection Potential decreases ($\Phi$ decreases).
    \item As vibration decreases ($f \to 0$), entropy decreases ($S \to 0$) and the Perfection Potential increases ($\Phi \to \infty$).
\end{itemize}
Black holes are the phase transducers effecting this transition. With the stopping of time ($t' \to 0$) and the nulling of frequency ($f' \to 0$) at the event horizon:
\begin{equation}
\lim_{f \to 0} S(f) = 0 \implies \Phi \to \infty
\end{equation}
This is the state of absolute order with zero entropy and infinite potential.

\section{Philosophical Interpretation: The Ontological Silence of Time}
If time is a phenomenon emerging with the vibration of strings,
then the stopping of time means the transformation of existence into silence.
At that point, the universe reaches the level of ``motionless perfection''.

This approach also agrees with Wheeler's ``It from Bit'' proposition:
being (``It'') is born from information (``Bit'').
In our model there is a cycle between these two states:
\[
\text{Energy (Vibration)} \rightleftharpoons \text{Information (Stationarity)}
\]

\section{Cosmological Cycle and Conclusion}
The universe is assumed to have a cyclic structure:
\begin{itemize}
\item \textbf{Beginning (Big Bang):} transformation of the stationary information phase ($\Phi \to \infty$) into energy ($f > 0$).
\item \textbf{Evolution:} the era in which strings vibrate, time flows, entropy increases ($S > 0$).
\item \textbf{Transformation:} black holes convert matter back into the stationary information phase ($f \to 0, S \to 0, \Phi \to \infty$).
\item \textbf{Rebirth:} information may turn back into energy, initiating a new universe.
\end{itemize}

This study defines black holes as phase transducers by bringing a string-based explanation to Einstein's time dilation equation and introducing the concept of Perfection Potential.
The slowing and stopping of time is not merely a relativistic effect, but the signature of the transition of strings from energy ($\Phi \to 0$) to information ($\Phi \to \infty$).

\begin{quote}
\textit{``When the vibration of the universe slows, time slows too;\\
when time stops, perfection begins.''}\\
--- Fahri Meral, Stationary Strings Theorem (2025)
\end{quote}

\section*{References}
\begin{enumerate}
\item Meral, F. (2025). \textit{Duran Sicimler Teoremi: Zamanın ve Entropinin Yokluğunda Mükemmel Düzenin Yeniden Doğuşu.}
\item Einstein, A. (1916). \textit{Die Grundlage der allgemeinen Relativitätstheorie}. Annalen der Physik, 49, 769–822.
\item Schwarzschild, K. (1916). \textit{\"Uber das Gravitationsfeld eines Massenpunktes nach der Einsteinschen Theorie}. Sitzungsberichte der Königlich Preussischen Akademie der Wissenschaften.
\item Hawking, S. W. (1974). \textit{Black hole explosions?} Nature, 248(5443), 30–31.
\item Hawking, S. W., \& Ellis, G. F. R. (1973). \textit{The Large Scale Structure of Space-Time}. Cambridge University Press.
\item Bekenstein, J. D. (1973). \textit{Black holes and entropy}. Physical Review D, 7(8), 2333–2346.
\item Maldacena, J. (1998). \textit{The Large N Limit of Superconformal Field Theories and Supergravity}. Advances in Theoretical and Mathematical Physics, 2, 231–252.
\item Maldacena, J., \& Susskind, L. (2013). \textit{Cool horizons for entangled black holes}. Fortschritte der Physik, 61(9), 781–811. (ER=EPR proposal)
\item Penrose, R. (2010). \textit{Cycles of Time: An Extraordinary New View of the Universe}. Bodley Head.
\item Polchinski, J. (1998). \textit{String Theory Vol. 1 \& 2}. Cambridge University Press.
\item Green, M. B., Schwarz, J. H., \& Witten, E. (1987). \textit{Superstring Theory}. Cambridge University Press.
\item Kaku, M. (1999). \textit{Introduction to Superstrings and M-Theory}. Springer.
\item Greene, B. (1999). \textit{The Elegant Universe: Superstrings, Hidden Dimensions, and the Quest for the Ultimate Theory}. W. W. Norton \& Company.
\item Susskind, L. (2008). \textit{The Black Hole War: My Battle with Stephen Hawking to Make the World Safe for Quantum Mechanics}. Little, Brown and Company.
\item Wheeler, J. A. (1990). \textit{Information, Physics, Quantum: The Search for Links}. In W. H. Zurek (Ed.), \textit{Complexity, Entropy, and the Physics of Information} (pp. 1–28). Addison-Wesley.
\item Page, D. N. (1993). \textit{Information in black hole radiation}. Physical Review Letters, 71(23), 3743–3746.
\item 't Hooft, G. (1993). \textit{Dimensional reduction in quantum gravity}. arXiv:gr-qc/9310026. (Holographic principle)
\item Srednicki, M. (1993). \textit{Entropy and area}. Physical Review Letters, 71(5), 666–669.
\item Bousso, R. (2002). \textit{The holographic principle}. Reviews of Modern Physics, 74(3), 825–874.
\item Preskill, J. (1992). \textit{Do black holes destroy information?} Proceedings of the International Symposium on Black Holes, Membranes, Wormholes and Superstrings.
\end{enumerate}

\end{document}
